\documentclass[11pt,a4paper]{article}
\usepackage[T1]{fontenc}
\usepackage[utf8]{inputenc}
\usepackage{amsmath,amssymb,xcolor}
\usepackage[margin=25mm]{geometry}
\usepackage[colorlinks=true,urlcolor=blue]{hyperref}
\definecolor{nexorablue}{HTML}{2459B5}
\setlength{\parindent}{0pt}
\setlength{\parskip}{7pt}
\setlength{\emergencystretch}{2em}
\renewcommand{\familydefault}{\sfdefault}
\begin{document}

{\small\bfseries\color{nexorablue} NEXORA CPGE}\par
{\small MP / NEX-STR-04 / Corrigé}\par
{\LARGE\bfseries Synchronisation de deux compteurs}\par
\vspace{0.25em}\hrule\vspace{0.6em}

\section*{Résolution}
On obtient $9k\equiv2\pmod{14}$. Comme $9\times11=99\equiv1\pmod{14}$, $k\equiv22\equiv8\pmod{14}$. Donc $t\equiv76\pmod{126}$.

\section*{Vérification}
Le reste de $76$ vaut $4$ modulo $9$ et $6$ modulo $14$. Comme les deux modules sont premiers entre eux, deux solutions diffèrent d’un multiple de $126$. Le plus petit entier positif est $76$.


\end{document}
